\documentclass[a4paper,12pt]{article}

\usepackage[utf8]{inputenc}
\usepackage[spanish]{babel}
\usepackage{geometry}
\geometry{margin=2.5cm}
\usepackage{graphicx}
\usepackage{amsmath}
\usepackage{amsfonts}
\usepackage{amssymb}
\usepackage{hyperref}
\usepackage{booktabs}
\usepackage{float}
\usepackage{subcaption}
\usepackage{listings}
\usepackage{xcolor}

% Configuración de listings para código Python
\lstset{
    language=Python,
    basicstyle=\ttfamily\small,
    keywordstyle=\color{blue},
    commentstyle=\color{gray},
    stringstyle=\color{red},
    breaklines=true,
    frame=single,
    showstringspaces=false
}

\title{Análisis Estadístico Avanzado del Rendimiento Estudiantil: \\
Aplicación de Inferencia Estadística y Manipulación de Datos}
\author{Alessandro Moises Britez}
\date{\today}

\begin{document}

\maketitle

\begin{abstract}
Este informe presenta un análisis estadístico completo del dataset de rendimiento estudiantil en matemáticas y portugués. Se aplicaron técnicas avanzadas de inferencia estadística incluyendo análisis de estimadores, demostración del Teorema del Límite Central, cálculo de intervalos de confianza, y manipulación eficiente de datos. Los resultados muestran que las calificaciones promedio en matemáticas (10.42 ± 4.58) y portugués (11.91 ± 3.93) difieren significativamente, con mayor variabilidad en matemáticas. Se demostró empíricamente la robustez de diferentes estimadores y la convergencia hacia distribuciones normales mediante simulaciones de bootstrap. El análisis confirma la eficiencia superior de pandas sobre métodos tradicionales de bucles para manipulación de datos.
\end{abstract}

\section{Introducción}

El rendimiento académico estudiantil es un tema de gran interés en el ámbito educativo, siendo fundamental comprender los factores que influyen en el éxito o fracaso escolar. Este trabajo presenta un análisis estadístico exhaustivo de un dataset que contiene información sobre el rendimiento de estudiantes portugueses en matemáticas y lengua portuguesa.

El objetivo principal de este estudio es aplicar técnicas avanzadas de inferencia estadística para analizar el rendimiento estudiantil, comparando diferentes estimadores estadísticos, demostrando propiedades fundamentales como el Teorema del Límite Central, y calculando intervalos de confianza para parámetros poblacionales. Adicionalmente, se evalúan técnicas eficientes de manipulación de datos, comparando metodologías tradicionales con enfoques modernos.

\subsection{Dataset}
El dataset utilizado proviene del repositorio UCI Machine Learning Repository y contiene información de 395 estudiantes de matemáticas y 649 estudiantes de portugués de dos escuelas secundarias portuguesas. Las variables incluyen calificaciones de tres períodos (G1, G2, G3), características demográficas, socioeconómicas y del entorno educativo.

\section{Metodología}

\subsection{Técnicas de Inferencia Estadística}

Para este análisis se implementaron múltiples técnicas estadísticas avanzadas:

\subsubsection{Análisis de Estimadores}
Se comparó la robustez y eficiencia de diferentes estimadores de tendencia central:
\begin{itemize}
    \item \textbf{Media aritmética}: $\bar{x} = \frac{1}{n}\sum_{i=1}^{n} x_i$
    \item \textbf{Mediana}: valor que divide la distribución en dos partes iguales
    \item \textbf{Desviación Absoluta Mediana (MAD)}: $MAD = \text{median}(|x_i - \text{median}(x)|)$
\end{itemize}

\subsubsection{Teorema del Límite Central}
Se realizaron simulaciones con diferentes tamaños de muestra ($n = 5, 10, 20, 50$) para demostrar que:
\[\frac{\bar{X}_n - \mu}{\sigma/\sqrt{n}} \xrightarrow{d} N(0,1) \text{ cuando } n \to \infty\]

\subsubsection{Intervalos de Confianza}
Se calcularon intervalos de confianza del 95\% para:

\paragraph{Media poblacional:}
\[IC_{1-\alpha}(\mu) = \bar{x} \pm t_{\alpha/2,n-1} \cdot \frac{s}{\sqrt{n}}\]

\paragraph{Varianza poblacional:}
\[IC_{1-\alpha}(\sigma^2) = \left[\frac{(n-1)s^2}{\chi^2_{\alpha/2,n-1}}, \frac{(n-1)s^2}{\chi^2_{1-\alpha/2,n-1}}\right]\]

\subsection{Bootstrap No Paramétrico}
Para validar los resultados teóricos, se implementó bootstrap con 10,000 remuestras para obtener intervalos de confianza empíricos y comparar con los métodos paramétricos tradicionales.

\subsection{Manipulación de Datos}
Se compararon dos enfoques para el análisis de datos:
\begin{itemize}
    \item \textbf{Método tradicional}: bucles iterativos (ineficiente)
    \item \textbf{Método moderno}: pandas/numpy vectorizado (eficiente)
\end{itemize}

\section{Resultados}

\subsection{Estadísticas Descriptivas}

Los resultados del análisis descriptivo muestran diferencias notables entre ambas materias:

\begin{table}[H]
\centering
\begin{tabular}{lcc}
\toprule
\textbf{Estadística} & \textbf{Matemáticas} & \textbf{Portugués} \\
\midrule
Número de estudiantes & 395 & 649 \\
Media (G3) & 10.42 & 11.91 \\
Desviación estándar & 4.58 & 3.93 \\
Rango & 0--20 & 0--20 \\
Tasa de aprobación (\%) & 67.1 & 85.4 \\
\bottomrule
\end{tabular}
\caption{Estadísticas descriptivas por materia}
\label{tab:descriptivas}
\end{table}

\subsection{Análisis de Estimadores}

El análisis de robustez reveló diferencias significativas entre estimadores:

\begin{table}[H]
\centering
\begin{tabular}{lcc}
\toprule
\textbf{Estimador} & \textbf{Valor Original} & \textbf{Con Outlier} \\
\midrule
Media & 10.415 & 11.570 (+11.1\%) \\
Mediana & 11.000 & 11.000 (+0.0\%) \\
MAD & 4.448 & 4.448 (+0.0\%) \\
\bottomrule
\end{tabular}
\caption{Robustez de estimadores ante outliers (Matemáticas)}
\label{tab:robustez}
\end{table}

\subsection{Eficiencia de Estimadores}

La comparación mediante bootstrap (1000 simulaciones) mostró:
\begin{itemize}
    \item \textbf{Varianza de medias bootstrap}: 0.1058
    \item \textbf{Varianza de medianas bootstrap}: 0.1377
    \item \textbf{Eficiencia relativa}: 1.301 (la media es 30\% más eficiente)
\end{itemize}

\subsection{Teorema del Límite Central}

Las simulaciones confirmaron la convergencia hacia distribuciones normales:

\begin{table}[H]
\centering
\begin{tabular}{lccc}
\toprule
\textbf{Tamaño muestra} & \textbf{Error estándar teórico} & \textbf{Error estándar observado} & \textbf{Error} \\
\midrule
n = 5 & 2.048 & 2.051 & 0.003 \\
n = 10 & 1.449 & 1.447 & 0.002 \\
n = 20 & 1.025 & 1.023 & 0.002 \\
n = 50 & 0.648 & 0.647 & 0.001 \\
\bottomrule
\end{tabular}
\caption{Convergencia del error estándar (Teorema del Límite Central)}
\label{tab:tlc}
\end{table}

\subsection{Intervalos de Confianza (95\%)}

\subsubsection{Matemáticas}
\begin{itemize}
    \item \textbf{IC Media}: (9.9621, 10.8686)
    \item \textbf{IC Varianza}: (18.4037, 23.7893)
    \item \textbf{IC Bootstrap Media}: (9.9547, 10.8725)
\end{itemize}

\subsubsection{Portugués}
\begin{itemize}
    \item \textbf{IC Media}: (11.6066, 12.2089)
    \item \textbf{IC Varianza}: (14.2847, 16.8456)
    \item \textbf{IC Bootstrap Media}: (11.6043, 12.2112)
\end{itemize}

La amplitud del IC para la media es menor en portugués (0.6023) que en matemáticas (0.9065), indicando mayor precisión en la estimación para portugués.

\subsection{Eficiencia Computacional}

La comparación de métodos de manipulación de datos mostró diferencias dramáticas:

\begin{table}[H]
\centering
\begin{tabular}{lcc}
\toprule
\textbf{Método} & \textbf{Tiempo (ms)} & \textbf{Mejora} \\
\midrule
Loop tradicional & 3.21 & - \\
Pandas vectorizado & 0.98 & 3.3x más rápido \\
\bottomrule
\end{tabular}
\caption{Comparación de eficiencia computacional}
\label{tab:eficiencia}
\end{table}

\subsection{Análisis de Variables Derivadas}

Se crearon nuevas variables para análisis más profundo:
\begin{itemize}
    \item \textbf{Progreso}: Diferencia G3--G1 (media: 0.29)
    \item \textbf{Balance tiempo}: Studytime--Freetime
    \item \textbf{Rendimiento categorizado}: Bajo/Medio/Alto basado en G3
\end{itemize}

\section{Conclusiones}

\subsection{Principales Hallazgos}

\begin{enumerate}
    \item \textbf{Diferencias entre materias}: El rendimiento en portugués es superior al de matemáticas, tanto en media como en tasa de aprobación (85.4\% vs 67.1\%).

    \item \textbf{Robustez de estimadores}: La mediana demostró ser completamente robusta ante outliers extremos, mientras que la media mostró sensibilidad significativa (+11.1\% de cambio).

    \item \textbf{Eficiencia estadística}: A pesar de su vulnerabilidad a outliers, la media es 30\% más eficiente que la mediana en términos de varianza muestral.

    \item \textbf{Validación del TLC}: Las simulaciones confirmaron empíricamente la convergencia hacia distribuciones normales, con errores menores a 0.003 en todos los casos.

    \item \textbf{Precisión de intervalos}: Los intervalos de confianza bootstrap coincidieron estrechamente con los métodos paramétricos, validando ambos enfoques.

    \item \textbf{Eficiencia computacional}: Los métodos vectorizados son 3.3 veces más rápidos que los bucles tradicionales, demostrando la importancia de técnicas modernas de programación.
\end{enumerate}

\subsection{Implicaciones Metodológicas}

Este estudio demuestra la importancia de:
\begin{itemize}
    \item Evaluar múltiples estimadores según el contexto y presencia de outliers
    \item Validar resultados teóricos mediante simulaciones
    \item Utilizar métodos no paramétricos como complemento a técnicas tradicionales
    \item Implementar código eficiente para análisis de grandes datasets
\end{itemize}

\subsection{Limitaciones y Trabajo Futuro}

Las limitaciones incluyen el análisis de un solo contexto educativo (Portugal) y la necesidad de considerar factores causales. Futuras investigaciones podrían explorar:
\begin{itemize}
    \item Modelos predictivos del rendimiento final (G3)
    \item Análisis longitudinal de la progresión G1 → G2 → G3
    \item Clustering para identificar perfiles de estudiantes
    \item Análisis causal de factores socioeconómicos
\end{itemize}

\section{Referencias}

\begin{enumerate}
    \item Cortez, P., Silva, A. (2008). Using Data Mining to Predict Secondary School Student Performance. EUROSIS.
    \item Casella, G., Berger, R. L. (2002). Statistical Inference, 2nd Edition. Duxbury Press.
    \item Efron, B., Tibshirani, R. J. (1994). An Introduction to the Bootstrap. Chapman \& Hall.
    \item McKinney, W. (2017). Python for Data Analysis, 2nd Edition. O'Reilly Media.
\end{enumerate}

\end{document}